\documentclass[11pt]{letter}
\usepackage[margin=1in]{geometry}
\usepackage{hyperref}
\usepackage{xcolor}

\signature{Ana Conesa\\on behalf of all authors}
\address{Institute for Integrative Systems Biology (I2SysBio)\\
Spanish National Research Council (CSIC)\\
Paterna 46980, Spain\\
ana.conesa@csic.es}

\date{\today}

\begin{document}

\begin{letter}{%
Dr.\ Marta del Olmo\\
Associate Editor, Interface\\
Nature Communications}

\opening{Dear Dr.\ del Olmo,}

Thank you for handling the review of our manuscript entitled
\emph{``TUSCO: benchmarking transcriptome reconstruction with endogenous single-isoform controls.''}
We are grateful to the editor and all five reviewers for their thorough
evaluation. We are pleased that Reviewers~\#1--3 and \#5 consider the
manuscript ready for publication, and we thank Reviewer~\#4 for their
additional minor comments, which have further improved the accuracy of
the text.

\bigskip
\textbf{Summary of revisions.}
Following Reviewer~\#4's comments, we have made four targeted corrections
in the manuscript, all highlighted in {\color{red}red font}:

\begin{enumerate}
  \item \textbf{Comment~1a} (line~272): Corrected the description of the
    MAJIQ-L metric to accurately state that it captures each method's
    junction-level sensitivity against a common short-read-derived splice
    graph, rather than reflecting sequencing depth.

  \item \textbf{Comment~1b} (lines~420--426): Rewrote the Discussion
    paragraph on short-read evaluation to clarify that (i)~MAJIQ-L
    provides a complementary junction-level assessment rather than
    exposing a short-read limitation, and (ii)~the affected transcripts
    are longer, multi-exon transcripts whose 5'-proximal junctions are
    harder to reach from the poly-A tail.

  \item \textbf{Comment~2} (lines~677--686): Added a Methods subsection
    detailing the MAJIQ-L pipeline configuration, specifying the build
    command, annotation version, junction scope (all junctions from all
    genes), and confirming that all parameters were set to their default
    values.

  \item \textbf{Comment~3} (lines~428--435): Added a Discussion paragraph
    explicitly acknowledging that the datasets evaluated vary in
    sequencing depth and that technology-level conclusions would require
    normalization strategies such as matching sequencing cost or total
    bases sequenced.
\end{enumerate}

\bigskip
\textbf{TUSCO tutorial fix (Comment~4).}
The tutorial error reported by Reviewer~\#4 stemmed from a mismatch in
data preparation: the annotation GTF was extracted with a narrower
flanking window than the accompanying genome FASTA, truncating
neighbouring genes at region edges. We have regenerated the tutorial
annotation GTF using the correct extraction window and also fixed a
related issue where the TUSCO report R~script crashed due to regex-based
gene ID matching. Both fixes have been submitted as pull request~\#588 to
the SQANTI3 repository
(\url{https://github.com/ConesaLab/SQANTI3/pull/588}) and verified on
both macOS and Linux.

\bigskip
\textbf{Additional materials.}
A detailed point-by-point response to all reviewer comments is attached
as a separate document.

\bigskip
\textbf{Compliance.}
We confirm that the revised manuscript complies with all Nature
Communications editorial policies. Data and code availability statements
are included in the manuscript. All computational code is publicly
available in the SQANTI3 GitHub repository and the TUSCO paper data
archive. The reporting summary checklist and the code and software
submission checklist are provided with this resubmission.

\closing{Yours sincerely,}

\end{letter}
\end{document}
